\chapter{11}

\section{Zürich (CH)}



Louis nickte rasch, hob den Rucksack hoch und bedankte sich.

Die Bremsen quietschten. Zwischen Leonards Anhalten und Louis’
Aussteigen riefen sie sich Worte der Verabschiedung zu. Winkten
einander. Leonard hupte zwei Mal kurz, scharf und laut. Schon sah Louis
dem mächtigen Truck nach. Der verschwand brummend und langsam in der
Kolonne zwischen den unzähligen andern Fahrzeugen.

Es war geschafft. Louis war als «Giovanni» durchgekommen, als «Gio-Gio»,
wie ihn Giovanni frischweg genannt hatte. Giovanni. Ihn würde er wohl
nie mehr sehen.

\sectionbreak

Zürich. Die Autos bretterten von allen Seiten an ihm vorbei.

Es roch nach alter Luft und staubigen Strassen. Louis lief der
Schweiss übers Gesicht. Er fühlte sich elend. Sein weiterer Weg lag völlig im
Dunkeln. Angst und Traurigkeit legten an Schärfe zu. Als hätten sie nur
darauf gewartet, bis er wieder mit sich alleine war. Listig und
erbarmungslos.

Die Mittagshitze hielt ihn mit Härte umfangen. Ein August der
Superlative, hiess es, pausenlos heiss, Tropennächte in Folge. Die
Menschen schwitzten sich teilweise geradezu in die Ohnmacht oder lagen
herum wie narkotisierte Fliegen.

Louis seufzte. Er musste sich anstrengen. Den Tag überstehen und
schliesslich später irgendwo einen Schlafplatz finden.

Er ging verlangsamt. Bei jedem Schritt wurde ihm unwohler. Übelkeit
stieg hoch. Er schwankte und kam sich vor wie ein ausgesetztes Kind.
Verlassen. Und auf sich selbst gestellt.

Es war seine Schuld. Er hatte es sich selbst angetan. Sich in eine
ausweglose Situation hineingeritten. Er war abgehauen. Es gab nichts
schön zu reden. Er war auf der Flucht. Er würde weiterhin Geschichten
zusammenlügen müssen. Und hatte keinen, dem er sich anvertrauen konnte.
Es schien ihm unmöglich, jemanden zu kontaktieren. Jedes seiner
Lebenszeichen konnte nachverfolgt werden. Man würde ihn aufspüren.

Giorgia? Seine Freunde? Marco? Maman? Seine Schwester?

Was würden sie alle denken, warum er verschwunden war. Noch keinen Tag
unterwegs und er fühlte sich bereits komplett ausgelaugt.
Und. Seine Gitarre fehlte.

Sobald er etwas ausserhalb war, sank er in sich zusammen. Was, wenn
Giorgia gar nicht mehr…Würde man einen wie ihn einen Mörder nennen?

Die Frage hängte sich unbeantwortet vor ihm auf. Daneben sah er
schlagartig seinen Vater stehen und erschrak.

Louis’ Übelkeit nahm zu. Vielleicht sollte er etwas essen.

Das Einkaufszentrum lag am Weg. Er durchschritt die einzelnen Geschäfte.

Er fand, was er brauchte. Als Lichtblick stellte sich der Kauf
der Trekkingschuhe zum halben Preis heraus. Die Lebensmittel und
Getränke mussten vorläufig reichen. Eine Taschenlampe mit
Ersatzbatterien als Sonderangebot legte er an der Kasse als letztes
hinzu.

Die Suche nach einer geeigneteren Brille war erfolglos geblieben.

\sectionbreak

Der Einkauf hatte ihn angestrengt. Wenigstens hatte sich niemand für ihn interessiert. 
Im Gegenteil, die Leute hier in Zürich waren strammen Schrittes unterwegs. Schauten weder links noch rechts.
Die meisten blickten auf ihr Handy oder sprachen vor sich hin. In ihren
Ohren steckten Kopfhörer. Auf den ersten Blick sahen sie aus, als wären
sie vor sich hin schwatzende Autisten. Einige Menschen rannten geradezu
an ihm vorbei. Als wären sie auf der Flucht. Wie er.

Womöglich stimmte, was Franco unlängst von einer Studie erzählt hatte. 
New York und Zürich seien die schnellsten Städte der Welt. Man
hatte die Geschwindigkeit untersucht, mit der sich die Menschen
bewegten. Es konnte ihm recht sein.

\sectionbreak

Während des Gehens riss Louis die Papiertüte eines Sandwiches auf.
Welchen Weg er auch immer einschlagen würde, es war einerlei. Louis
versuchte, sich auf den Geschmack beim Kauen zu konzentrieren. Seine
Gedanken kamen ihm vor wie ungebetene Gäste, die lärmend vor seiner Tür
Schlange standen, um sich Einlass zu erzwingen. Es war aussichtslos, sie
zu ignorieren. Er seufzte halblaut vor sich hin.

Unerwartet tat sich in der Ferne eine schattenspendende Allee auf. Louis
beschleunigte das Tempo. Und wie auf ein Signal fiel ihm ein, wie er
sich kürzlich in seiner Stimmung von Traurigkeit \emph{«Ocean»} von
John Butler angehört und sich darin aufgehoben gefühlt hatte.

Er hatte dagesessen, nachdem sich Giorgia das erste Mal deutlich gegen
ihre Beziehung ausgesprochen hatte. Wollte nur noch alleine sein. Er
hatte seine Gitarre genommen, sich das Stück zig Mal angehört, die
Harmonien bestimmt und es nachgespielt. Erst holprig, wiederholend.
Eines Tages hatte es geklungen.

Er hatte alleine für sich gespielt. Immer wieder, wenn er einbrach.
Konnte sich in das virtuose Spiel hineinlegen und darin versinken.
Entschwinden und kurzzeitig hinaus finden aus seinen düsteren
Stimmungen. Das Spiel schien seine Verlorenheit aufzufangen. Vielleicht
ersetzte es auch seine Tränen.

Louis’ Schritte blieben schleppend. Die Umgebung dünnte sich aus.
Weniger Verkehr, weniger Häuser, weniger Menschen.

Wieder tränten seine Augen. Es lag wohl an der Hitze. Das penetrante
Gefühl der Angst schlich sich erneut an. Ein einzelner Griff hätte ihm
alles erspart, \emph{che stronzo sono}\footnote{…was bin ich für ein Arschloch.}! Augenblicklich begann «Ocean»
in ihm zu spielen. Seine Finger bewegten sich. Mit der rechten Hand
imitierte er rhythmisches Schlagen auf Hals und Klangkörper seiner
\emph{Taylor}. Die versessen trainierte Golpe-Technik hatte er sich bei
einem Flamencogitarristen abgeschaut. Es gefiel ihm, sein eigener
Rhythmus zu sein.

Er legte summend eine Grundmelodie über die Klangfolge. Verhalten erst
und leise. Dann war ihm, als fliesse seine Trauer ungebrochen in den
geöffneten Tonkanal.

Die Angst verlor an Dominanz.

\qrblock{John Butler \emph{«Ocean»}}{media/QR-Codes/013_Ocean_John_Butler_Spotify.png}

Unerwartet hörte er Grandpères Worte in sich.

\sectionbreak
\indentedblock{Damals. Fribourg, sieben.

«Angst hat gute und schlechte Seiten. Wenn deine Angst dich nervt,
Louis, kannst du sie reinlegen.»

Sie sitzen auf einer Bank an der Strasse, Grandpère und er.

Grandpère zwinkert ihm zu. Louis schaut zu ihm hoch und schmunzelt
erwartungsvoll. Was mit Tricks zusammenhängt, fasziniert ihn.

«Also, stell dir vor, in deinem Gehirn im Kopf gibt es viele kleine
Gefühlsschubladen. In einer wohnt die Aufregung, in einer anderen die
Freude und wieder in einer andern die Angst, und all deine anderen
Gefühle. Dann gibt es noch dutzende weitere Schubladen. Da wohnen Dinge,
die wir gelernt haben. Das Lesen und Rechnen und so weiter…alles, was
du gerade in der Schule so lernst. In einer wohnt das Singen. Das Gehirn
ist so gemacht, dass manchmal Schubladen gleichzeitig offen sind…»

«Gleichzeitig, wie denn?» fragt Louis gespannt.

«Hm, warte.»

Grandpère denkt nach. Louis lutscht weiter an seinem Lolly. Dabei
fixiert er erwartungsvoll Grandpères Gesicht.

«Also, stell dir vor, du guckst dir einen Film an. Da sind zum Beispiel
Schublade Schauen, Hören, Freude, Aufregung, vielleicht Angst offen, es
kommt auf den Film an.»

Louis nickt motiviert. Er kann Grandpère folgen.

«Es gibt Schubladen, die sich nicht gleichzeitig öffnen lassen. Das ist
gut so, und wenn du Angst hast, sogar grossartig, fast schon magisch.»

«Und?»

Louis hebt die Augenbrauen. Wackelt mit dem Kopf. Er ist ungeduldig und
kann es kaum erwarten, dass Grandpère auf den Punkt kommt.

«Wenn die Singschublade offen ist, ist die Angstschublade blockiert.
Auch umgekehrt, wenn die Angstschublade offen ist, ist die Singschublade
zu. Einfacher, Louis, du kannst nicht gleichzeitig singen und Angst
haben.»

Dann lächelt Grandpère ihn gütig an. Er streichelt ihm übers Haar.

«Wenn du also Angst hast, zum Beispiel vor grösseren Jungs auf dem
Schulweg oder wenn du im Dunkeln alleine draussen bist, oder sonst wann,
Louis, dann fang an zu singen, auch summen geht und auch leise singen
klappt.»

«Und die Angst ist weg, Grandpère?»

«Hm,» nickt dieser überzeugt, «singen und gleichzeitig Angst haben geht
nicht. Das Gehirn kann das nicht.»

Es ist verrückt. Es klappt.}

\sectionbreak

Nach längerem Gehen erreichte er ein grosses Gebäude im Rohbau am Rande
der Stadt. In weitem Abstand dazu lagen Mehrfamilienhäuser.
